\scanpage{80}
\mspage{24}

Comme $S'\unsure{\text{---}|\to}S$ est surjectif, cela montre que les $f_s:X_s\unsure{\text{---}|\to}Y_s$ sont dans $W$, qed.

Ainsi, (2) sera conséquence de

\textcircled{$2_0$} Soit $p:\deleted{\unsure{X'}}Y'\to\deleted{\unsure{X}}Y$ une Hot-fibration (ou \inserted{même} un foncteur parfait, si ça peut marcher \ldots). Soit $f:X\to Y$ une $S$-flèche (dans $\mathrm{Cat}_{/S}$), \unsure{déduis|déduisons|obtenu} $f':X'\unsure{\text{---}|\to}Y'$ par chgt. de base. Si $f\in W$, $f'\in W'$\,:
\diagram{\begin{tikzcd}[ampersand replacement=\&,column sep=small,row sep=small]
X'\arrow[r,"q"]\arrow[d,"f'"']\&X\arrow[d,"f"]\\
Y'\arrow[r,"p"']\&Y
\arrow[from=1-1,to=2-2,phantom,"\mathrm{cart}" description]
\end{tikzcd}}
\[ f\in W\ \text{et}\ p\ \text{parfait}\ \deleted{\unsure{\text{sur}}}\ \Longrightarrow f'\in W \]

En tous cas, il nous faut absolument \textcircled{2} pour pouvoir conclure que $C$ est compatible avec les localiseurs.

Considérons $rC:\mathrm{Cat}_{/S}\unsure{\longrightarrow|\text{---}}\mathrm{Cat}_{/S}$ on a une flèche
\[
\left\{\begin{aligned}
&\mathrm{id}_{\mathrm{Cat}_{/S}}\xrightarrow{\deleted{\unsure{i}}j}rC\\
&\deleted{\unsure{i}}j_X:X\longrightarrow CX\qquad \substack{\text{la flèche canonique}\\ i_{X/S}\ (\text{p.~21})},
\end{aligned}\right.
\]
et $j_X$ est tjrs dans $W$.

Considérons $Cr:\underline{\mathrm{Parf}}\,S\unsure{\text{---}|\to}\underline{\mathrm{Parf}}\,S$, i.e appliquons la construction $C$ à un $X$ parfait sur $S$. Alors on a encore
\[
\left\{\begin{aligned}
&\mathrm{id}_{\underline{\mathrm{Parf}}\,S}\xrightarrow{j'}Cr\\
&j'_X=j_{\unsure{X|r}}(X):\overset{\deleted{\unsure{r(X)}}}{X=r(X)}\xrightarrow{j_X}C(X).
\end{aligned}\right.
\]

Ledit morphisme est dans $W$, donc aussi une $W$-équiv fibre par fibre, puisque \unsure{sources|source} et \unsure{buts|but} \deleted{sont} parfaits, donc des Hot-fibrés sur $S$.

% archive sheet number: 79
